
\section{How We Receive What We Need}

Installment 2: “Folkebladet,” March 01, 1916

The question, then, is this: how can we come to share in these blessings, so glorious and indispensable as they are for us?

There may perhaps be one person or another who, without further thought, consoles himself with the belief that we possess them, and therefore need not concern ourselves with this matter, either for ourselves or for others.

I do not think, however, that very many think in this way. Most among us must surely acknowledge that we are unspeakably poor in spiritual life and in everything else that we so deeply need, both in our personal life of faith and in our congregational life.

We therefore feel the importance of the question: how can we receive what we need? It is perhaps not so easy to answer this question.

I believe, however, that many agree with me when, to begin with, I answer: We never receive these blessings and gifts against our will, nor do we receive them without doing something in order to receive them.

Yes, but what can we do? Perhaps not much; but we can do something, provided only that we are willing. We can, first of all, acknowledge our spiritual poverty and our lukewarmness. But this is by no means an easy road to walk. It costs humiliation and suffering, and from that we shrink.

I do not know for certain, but it may be that there has been some inclination among us to settle into complacency both with our own spiritual condition and with that of our congregations. Yes, that is dangerous.

Paul was never satisfied, either with his own relation to God or with that of the congregations. Always more, always farther onward: “I press on to lay hold of it”—“I hasten toward the goal.”

And he exhorts the congregations to which he writes his letters to do likewise: “You also, brothers, become imitators of me!” “Therefore I bow my knees before the Father … that he … may grant you to be strengthened with power through his Spirit in your inner being.”

Yes, if this deep acknowledgment of our powerlessness and our half-hearted Christianity could seize us, fill us! Then there would be cries for help, and then we would receive help.

To his first congregation Jesus said: “You will receive power when the Holy Spirit comes upon you, and you will be my witnesses both in Jerusalem and to the ends of the earth.”

And of this congregation it is written: “They continued with one accord in prayer.” And they continued in this prayer with one accord until, on the day of Pentecost, they received the power from on high.

He who alone can give us what we most need, also in our Free Church work—he has said to us as well: “Ask, and you will receive!” We can, if we are willing, acknowledge our helplessness and need and lament them before the Lord. We can tell him at least something about what is lacking—not first and foremost in others, but in ourselves, both individually and as congregations.

Certainly he knows it well, he who says, “I know your works,” without our telling him; but he has nevertheless laid down in his word the rule that spiritual blessing reaches neither the individual nor the congregation without earnest prayer. And here, no doubt, lies our great loss: We pray so little, both in the private chamber and in the congregational gathering. I know well that believing people can never cease praying altogether; but I also know that the heart can grow cold, the knees stiff, and prayer many a time become only a prayer for ourselves.

Brothers and sisters, let us attend to this matter. For it is so important.

The kingdom of God that Jesus brought was not intended to be merely a Christianity of individuals. It was intended to be a Christianity of fellowship, in the sense in which we confess in the Third Article: “I believe in the communion of saints.”

It is certainly true that it is as individuals that we come to share in salvation and enter into union with God; but it is not intended that each believer should live his life in God for himself alone, but rather that he should live this life in the blessed consciousness of belonging among the Lord’s great company of children.

According to the New Testament, there can be no mistaking that it is God’s will and purpose that his people should live their life in God as a life in fellowship. Therefore so much is said about “fellowship.”

And Jesus prays for his own that “they may be one” with almost greater earnestness than for anything else.

We may therefore conclude that, for him, this matter has overwhelmingly great importance. Yes, he himself says that he prays for this “so that the world may believe that you have sent me.”

The progress and spread of Christianity depend upon this.

It is therefore important to think of the congregation and the congregations when the question is how we are to receive the gifts of grace without which the Free Church cannot be built. And when it concerns the manner in which these gifts may be obtained, the way is not simply the same for the individual as for the congregation. For the individual, perhaps one of the most important things required in order to receive these gifts is that solitude be sought. For the congregation, on the other hand, what is required, as already pointed out, is that fellowship be sought—the communion of saints.

Even after the outpouring of the Spirit, it is said of the first Christians: “They held faithfully to the apostles’ teaching and to the fellowship, to the breaking of bread and to the prayers.” Acts 2:42.

They therefore found it necessary to gather with one another in such assemblies, where love for one another and love for the Lord and his cause were the sole command impelling them to gather.

Concerning this matter, the late Professor Sverdrup says in Collected Writings, volume 3, page 307:

\begin{quote}
If free, independent congregations are not to lose something of the power of cooperation, spiritual connection is needed. Free cooperation in everything that can serve the advancement of God’s kingdom at home and abroad is necessary. If it is to take place, then we must understand one another and know one another. We must exchange thoughts and spirit and thus have the bonds of love tied in the freedom of the Spirit. As far as I can understand, genuine Free Church work will be almost impossible under our conditions without a great many free meetings and gatherings.
\end{quote}

As may be seen from these statements by Professor Sverdrup, he placed great weight upon the people of the Free Church coming together for free meetings.

And it is a well-known fact among us that quite a number of such larger or smaller meetings are also held, at which it is desired that brothers and sisters from several congregations should gather. This is both necessary and right; for we shall scarcely be supplied in any other way with that power from on high which we must have, if we are to be able to fulfill our task, than by gathering for such meetings.

But I, and perhaps many with me, have observed with concern that the felt need to come together for annual meetings, mission meetings, district meetings, and so forth seems to be declining to an alarming degree. What can be the cause? Could it perhaps be the same cause for which the congregation and the angel of the congregation in Laodicea had to hear these solemn words: “Because you say, ‘I am rich and have abundance and lack nothing,’ and you do not know that you are wretched and pitiable and poor and blind and naked, I therefore counsel you to buy from me—”?

Oh, let us follow his counsel: “buy from him.”
